\section*{第 \ref{ch:g2:measuring-time} 章——测量时间}

\begin{solution}{exo:g2:measuring-time:1}
起床、吃早餐、下午玩耍、吃晚饭、睡觉。
\end{solution}

\begin{solution}{exo:g2:measuring-time:2}
\omterm{def:g2:measuring-time:duration}{时刻}；\omterm{def:g2:measuring-time:duration}{时长}；\omterm{def:g2:measuring-time:duration}{时刻}；\omterm{def:g2:measuring-time:duration}{时长}。
\end{solution}

\begin{solution}{exo:g2:measuring-time:3}
蓝色蜡烛烧得更久。两根蜡烛同时点燃，
另一根燃尽时仍在燃烧的那根，就是烧得更久的——
这正是公平赛跑的规则。
\end{solution}

\begin{solution}{exo:g2:measuring-time:4}
每次翻转后，沙子穿过狭窄的腰部，所用\omterm{def:g2:measuring-time:duration}{时长}都\emph{相同}。
一个从不改变的计时器，
对每个人都是公平的。
\end{solution}

\begin{solution}{exo:g2:measuring-time:5}
我们对时间的感觉是一个不准的钟：无聊的等待慢如爬行，快乐的时光飞一般流逝。
不论本有什么感觉，手表都稳定地计时，所以等待的时间其实
比动画片还短——它只是\emph{感觉}更长。
\end{solution}

\begin{solution}{exo:g2:measuring-time:6}
答案不一；一般说出名字约需几秒，
系鞋带约需 $20$ 秒。数数时要保持平静、稳定。
\end{solution}

\begin{solution}{exo:g2:measuring-time:7}
打喷嚏：秒。整夜睡眠：\omterm{ex:g2:measuring-time:clock}{小时}。系鞋带：
秒。上学的早晨：\omterm{ex:g2:measuring-time:clock}{小时}。
\end{solution}

\begin{solution}{exo:g2:measuring-time:8}
五点整时，时针指向 $5$（分针竖直向上）；
五点半时，时针指向\emph{$5$ 与 $6$ 的正中间}
（分针竖直向下）。
\end{solution}

\begin{solution}{exo:g2:measuring-time:9}
用同一个沙漏做比较：洗澡时，开始计时，
数一数洗澡持续了沙漏的几个完整回合；明天，
对淋浴做同样的事。记下两个数（比如“洗澡：$4$
次，淋浴：$2$ 次”）：次数越多，时间越长；任何人
拿同一个沙漏都能检验。
\end{solution}

\begin{solution}{exo:g2:measuring-time:10}
汤姆更接近正确。只有在平静、稳定的节奏下，数秒才真实可靠——
每个数字大约对应一次平静的心跳，或者每数一下念一声“一只小鳄
鱼”。飞快地数数字，
会让每一秒变得太短。
\end{solution}
